\chapter{Transforming words in parallel}
\MARSChapterHeading{Transforming words in parallel}
\label{revision:chacha}

A vector instruction is useful when several calculations have the same shape.
The difficulty is often arranging the data so that the lanes describe independent
calculations. ChaCha20 gives us a small example that we can first understand
with four scalar values, then reorganize without changing the answer.

This chapter develops the complete program in
\texttt{revision/labs/files/chacha20.asm}. It produces a stream of bytes from a
key, nonce and block counter, and combines those bytes with a counted message.
The word transformation is defined by
\href{https://www.rfc-editor.org/rfc/rfc8439.html}{RFC 8439, sections 2.1--2.4}.
Our function interfaces and failure rules are specified separately in package
\texttt{026-chacha20}. Read the implementation as an exercise in precise word
arithmetic, data arrangement and ownership. Designing a cryptographic algorithm
is a different task from implementing its specified operations.

We will build the calculation from the inside out. First trace four words. Then
identify which four words each call selects from a block. Only after that will
we arrange four simultaneous calls in MSA lanes. The byte wrapper comes last:
it depends on a correct block calculation, but has its own bounds and state.

\section{Make the arithmetic width explicit}

A ChaCha word has 32 bits, even though a MIPS64 scalar register has 64. Addition
keeps the low 32 result bits. For example, adding hexadecimal
\texttt{ffffffff} and \texttt{00000002} produces word \texttt{00000001}.
There is no overflow error to report. The discarded carry is part of the
specified arithmetic.

The worked scalar code uses \texttt{addu} for this operation. Its suffix means
that signed overflow does not trap; the instruction still performs a word
addition. A MIPS64 word result is sign-extended into the register's upper half.
That upper half is not another part of the ChaCha word. When tracing this code,
compare eight hexadecimal digits for each word. Address arithmetic continues to
use full-width operations such as \texttt{daddiu}; narrowing an address would
be an unrelated error.

The other operations are exclusive OR and rotation. XOR changes each bit where
its two operands differ. A rotation moves bits around a word and puts the bits
that leave one end back at the other. Rotating \texttt{80000001} left by one
produces \texttt{00000003}. A left shift would instead discard the high bit and
produce \texttt{00000002}.

The scalar code expresses a left rotation by $r$ as a right rotation by
$32-r$. Thus \texttt{rotr} by 20 implements a left rotation by 12. These are
word rotations. Replacing them with doubleword rotations would change the
algorithm by circulating bits through the register's upper half.

\section{Follow the four dependencies in a quarter-round}

A \emph{quarter-round} updates four words named $a$, $b$, $c$ and $d$.
Each row below finishes before the next row starts. Every addition is modulo
$2^{32}$, and $\operatorname{rotl}$ rotates a 32-bit word to the left.

\begin{center}
  \begin{tabular}{lll}
    \textbf{Add} & \textbf{Combine}            & \textbf{Rotate}                   \\
    $a\gets a+b$ & $d\gets d\mathbin{\oplus}a$ & $d\gets\operatorname{rotl}(d,16)$ \\
    $c\gets c+d$ & $b\gets b\mathbin{\oplus}c$ & $b\gets\operatorname{rotl}(b,12)$ \\
    $a\gets a+b$ & $d\gets d\mathbin{\oplus}a$ & $d\gets\operatorname{rotl}(d,8)$  \\
    $c\gets c+d$ & $b\gets b\mathbin{\oplus}c$ & $b\gets\operatorname{rotl}(b,7)$
  \end{tabular}
\end{center}

\MARSCodeLabel{The first dependency group, with a word rotation}
\begin{minted}{mars-idea}
    addu $t0, $t0, $t1
    xor $t3, $t3, $t0
    rotr $t3, $t3, 16
\end{minted}

The first XOR uses the newly computed $a$. The second addition uses the newly
rotated $d$. These dependencies explain why the four lines cannot simply run
in parallel. Keeping the original operands in a second set of registers and
using them throughout would compute a different transformation.

The supplied caller starts with words \texttt{11111111}, \texttt{01020304},
\texttt{9b8d6f43} and \texttt{01234567}. After the first addition, $a$ is
\texttt{12131415}. XOR with the original $d$ gives \texttt{13305172}; rotating
that result left by 16 gives \texttt{51721330}. Stop at
\texttt{after\_qr\_group1} to inspect those first changes. At that boundary,
$b$ and $c$ still contain their input words.

After all four groups, the selected words must be \texttt{ea2a92f4},
\texttt{cb1cf8ce}, \texttt{4581472e} and \texttt{5881c4bb}, respectively.
A wrong final word tells us less than the first wrong intermediate word.
If the first group is already wrong, inspect addition width, XOR operands and
rotation direction before investigating any block loop.

\texttt{chacha\_quarter} receives four addresses in \texttt{a0} through
\texttt{a3}, not four immediate word values. It loads the selected words,
computes in temporary registers and stores to the same four addresses.
The caller deliberately spaces those words 20 bytes apart. A function that
loads the first address and assumes the other words are adjacent violates its
interface even if it happens to pass a contiguous test. The gap bytes are not
scratch space.

\begin{figure}[!ht]
  \centering
  \begin{tikzpicture}[alt={Actual light-theme native Memory view after the quarter-round. An outline encloses the first four result bytes f4 92 2a ea. The following twelve bytes remain a5; the other selected words are not pictured.}]
    \node[anchor=north west,inner sep=0,font=\small\bfseries] (title) at (0,0) {A word result and the untouched gap};
    \node[anchor=north west,inner sep=0] (a) at ([yshift=-2mm]title.south west) {\MARSRevisionCrop{revision/captures/chacha/chacha-quarter.png}{0.99\linewidth}{0.03125}{0.33888889}{0.58203125}{0.60888889}};
    \begin{scope}[shift={(a.south west)},x={($(a.south east)-(a.south west)$)},y={($(a.north west)-(a.south west)$)}]
      \draw[black,line width=.7pt] (.20,.01) rectangle (.39,.467);
    \end{scope}
  \end{tikzpicture}
  \caption{The first selected word is serialized as f4 92 2a ea, or word ea2a92f4. The next twelve pictured bytes remain a5. The other three selected words are outside this crop and are verified by the quarter-round checks.}
  \label{fig:revision-chacha-native-quarter}
\end{figure}


\section{Select positions before counting rounds}

A block contains sixteen words. Number them in row order:

\begin{center}
  \begin{tabular}{rrrr}
    0  & 1  & 2  & 3  \\
    4  & 5  & 6  & 7  \\
    8  & 9  & 10 & 11 \\
    12 & 13 & 14 & 15
  \end{tabular}
\end{center}

These numbers identify positions, not the values currently stored there.
Word position $i$ has byte offset $4i$. A column round applies one quarter-round
to each of the following selections, with the listed order giving $a,b,c,d$:
$(0,4,8,12)$, $(1,5,9,13)$, $(2,6,10,14)$ and $(3,7,11,15)$.
Those four selections do not overlap, so their quarter-rounds are independent.

A diagonal round follows, selecting $(0,5,10,15)$, $(1,6,11,12)$,
$(2,7,8,13)$ and $(3,4,9,14)$. These selections use the words left by the
column round. Together, the column round and diagonal round form one
\emph{double-round}. Ten repetitions give the twenty rounds in ChaCha20.
There are eight quarter-round calls per repetition and eighty per block.

In the scalar implementation, the destination initially receives a copy of the
sixteen input words. It is then the working block. The input remains unchanged.
The scalar routine saves its input pointer, output pointer and repetition count
across nested quarter-round calls. A temporary register cannot retain that
count merely because one particular version of the helper leaves it alone.

After the tenth double-round, the function adds each original input word to its
corresponding working word. This final addition is called \emph{feed-forward}.
The original state must still be available at this point. Adding the working
block to itself, or adding an already modified input block, produces the wrong
output despite a correct round loop.

Finally, serialize the sixteen sums in little-endian order. The worked
simulator uses little-endian memory, so word stores produce the required byte
order. A displayed word \texttt{e4e7f110} occupies memory as bytes
\texttt{10 f1 e7 e4}. A byte-oriented reference and a word-oriented debugger can
therefore show the same result in different-looking groups. Compare like units
before changing the calculation.

\begin{figure}[!ht]
\centering
\begin{tikzpicture}[alt={Actual light-theme native Memory view after the scalar block function. The first byte row starts 10 f1 e7 e4 d1 3b 59 15 and the next row starts c7 d1 f4 c7. An outline identifies the first four bytes.}]
\node[anchor=north west,inner sep=0,font=\small\bfseries] (title) at (0,0) {Words and bytes describe the same output};
\node[anchor=north west,inner sep=0] (a) at ([yshift=-2mm]title.south west) {\MARSRevisionCrop{revision/captures/chacha/chacha-block.png}{0.99\linewidth}{0.03125}{0.31555556}{0.58203125}{0.60888889}};
\begin{scope}[shift={(a.south west)},x={($(a.south east)-(a.south west)$)},y={($(a.north west)-(a.south west)$)}]
\draw[black,line width=.7pt] (.20,.3088) rectangle (.39,.6176);
\end{scope}
\end{tikzpicture}
\caption{The first 32 generated bytes are shown. The first four, 10 f1 e7 e4, encode word e4e7f110 in little-endian order. The outline identifies that word; the full 64-byte output is checked separately.}
\label{fig:revision-chacha-native-block}
\end{figure}


\section{Turn independent selections into vector lanes}

The MSA version loads one row into each of \texttt{w0} through \texttt{w3}.
Consider lane zero across those registers: it holds positions 0, 4, 8 and 12.
Lane one holds 1, 5, 9 and 13, and so on. A word addition between \texttt{w0}
and \texttt{w1} now performs the first $a+b$ operation for all four columns.
Each lane is an independent quarter-round; no operation adds one lane to its
neighbor.

\texttt{addv.w} performs four wrapping word additions. \texttt{xor.v} combines
the corresponding bits. A rotation is formed from word shifts and OR: shift
left by $r$, shift the same unmodified value right logically by $32-r$, then
combine the results. The two shifts need separate temporaries. Shifting an
already shifted value for the second half would lose the bits that should wrap.
The worked allocation uses \texttt{w8/w9} for those temporaries and retains
original rows in \texttt{w4} through \texttt{w7} for feed-forward.

After a column round, the lane arrangements must change before the diagonal
round. Keep row zero in place. Select row one's source lanes in order
$1,2,3,0$, row two's in order $2,3,0,1$, and row three's in order $3,0,1,2$.
Across lane zero we now have positions $(0,5,10,15)$, exactly the first diagonal.
Across lane one we have $(1,6,11,12)$. This is why a row selection can replace
four sets of scalar address calculations.

\begin{center}
  \begin{tabular}{lrrr}
    \textbf{Row} & \textbf{Select immediate} & \textbf{Inverse immediate} \\
    \texttt{w1}  & 57                        & 147                        \\
    \texttt{w2}  & 78                        & 78                         \\
    \texttt{w3}  & 147                       & 57
  \end{tabular}
\end{center}

\MARSCodeLabel{Arrange diagonal inputs without changing row zero}
\begin{minted}{mars-idea}
    shf.w $w1, $w1, 57
    shf.w $w2, $w2, 78
    shf.w $w3, $w3, 147
\end{minted}

Each two-bit field of a \texttt{shf.w} immediate names the source lane for one
destination lane, starting with destination lane zero in the low field.
For 57, the fields read from low to high are $1,2,3,0$. Decode the fields;
reading the binary literal from its printed left edge reverses their meaning.

After the diagonal arithmetic, apply the inverse selections. The next column
round, the original-row feed-forward and the final stores all expect the normal
row arrangement. Omitting the inverse does not just change the display: it
changes which words participate in the next repetition. A useful diagnostic is
to compare scalar and vector working states after one complete double-round,
when both should be back in the same arrangement.

The illustrated vector function clobbers \texttt{w0} through \texttt{w9},
including their floating-point register aliases. Its caller must respect that
published allocation. The completed function checks accept another correct
allocation; guided intermediate checks necessarily describe this particular one.

\section{Generate bytes without consuming the message early}

For the IETF stream interface, the initial state's first four words are fixed
constants. Eight words hold the 32-byte key. Word 12 is a 32-bit block counter,
and words 13--15 hold a 12-byte nonce. A nonce identifies a stream under a key;
it is an input field, not another portion of the counter.

The block function accepts any sixteen input words, which makes its arithmetic
independently testable. The stream wrapper instead receives a valid IETF state.
The wrapper copies that state into its own frame and increments only the private
counter between required blocks. The caller's state remains unchanged.

A generated block provides 64 \emph{keystream} bytes. For each selected message
byte, output is the source byte XOR the corresponding keystream byte. Applying
the same keystream again recovers the source, because $(x\mathbin{\oplus}k)
  \mathbin{\oplus}k=x$. Encryption and decryption therefore use the same byte
operation and require the same starting stream position.

The byte loop loads both operands before storing the result. This ordering
supports exact in-place transformation: source and destination can identify
the same byte. Arbitrary partial overlap is outside this exercise's interface.
Moving the destination one byte forward could overwrite a source byte needed
by the next iteration.

For a 65-byte message, generate one block and transform 64 bytes, then generate
another block and transform one byte. Generating a second block does not permit
reading another 64 message bytes. Unused keystream bytes are discarded by this
interface. Every load from the message and every store to its destination must
remain within the declared count, even though the private block buffer is larger.

\section{Reject an impossible request before writing}

\texttt{chacha\_xor} accepts an unsigned count from zero through 4096.
Check that full-width count first. An excessive count returns status one
without accessing any supplied pointer. A zero count returns success with the
same no-access promise. Only then may a positive request read its state.
Checking a narrowed word would mistakenly accept some large 64-bit counts.

For positive count $n$, let $B=\lfloor(n+63)/64\rfloor$. If the initial
counter is $c$, the final block uses counter $c+B-1$. Compute that expression
in a wider unsigned representation and reject a result above
\texttt{ffffffff}. The bound of 4096 makes the arithmetic small enough to fit
comfortably in a 64-bit register. Load $c$ with \texttt{lwu}; a signed word
load would give the wrong widened value when its top bit is set.

\begin{center}
  \begin{tabular}{lrll}
    \textbf{Initial counter} & \textbf{Bytes} & \textbf{Blocks} & \textbf{Result} \\
    \texttt{ffffffff}        & 1              & 1               & success         \\
    \texttt{ffffffff}        & 64             & 1               & success         \\
    \texttt{ffffffff}        & 65             & 2               & status two      \\
    \texttt{fffffffe}        & 65             & 2               & success
  \end{tabular}
\end{center}

The largest counter is usable. What fails is a request for a later block.
Perform this check before the first destination store. Rejecting only when the
loop eventually wraps would publish a prefix and then fail, violating the
whole-request promise. Nor may a counter increment carry into the nonce: those
bytes describe a different input field.

Return after the last required fragment, before incrementing the counter again.
This both avoids unnecessary work and permits the last available block.
Restore full saved registers and the stack pointer on every allocated-frame
exit. The worked wrapper's 192-byte frame and scalar block's 48-byte frame
preserve a caller's alignment through their nested calls.

\section{Distinguish a restart from a continuation}

The function retains no hidden stream position. Two one-byte calls with the
same state both use byte zero of the same block. They do not consume successive
bytes. To continue after a partial block, another interface would need to retain
its unused keystream and byte position, or explicitly request an offset.
Updating only the block counter after a one-byte call skips the other 63 bytes.

This also explains why changing a message while keeping its key, nonce and
counter repeats the same keystream. For two such outputs, XOR cancels that
keystream and exposes the XOR of the two messages. The exercise uses fixed
inputs for reproducible inspection; those fixtures are not a nonce-allocation
scheme.

Finally, a successful round trip proves neither authenticity nor complete
correctness. Two matching implementation mistakes can cancel during decryption.
Changing a ciphertext bit changes the corresponding recovered bit without
making this stream function report a failure. The authenticated construction in
RFC 8439 adds a separate authenticator; it is outside this program.
Use independent known outputs, arbitrary raw block inputs, exact memory bounds
and repeated-call tests to evaluate the implementation we actually wrote.

Before the lab, identify the first incorrect state caused by a doubleword
rotation, an omitted inverse shuffle, a signed counter load, a full-block read
for a one-byte tail, or a hidden counter. Name the violated promise before
changing instructions.
